% Part: second-order-logic
% Chapter: metatheory
% Section: introduction

\documentclass[../../../include/open-logic-section]{subfiles}

\begin{document}

\olfileid{sol}{met}{int}

\olsection{పరిచయం}

మొదటిస్థాయి తర్కానికి అనేక చక్కని ధర్మాలు ఉన్నాయి. అది నిర్ణయించదగినది
కాదని మనకు తెలుసు; కానీ కనీసం అది స్వీకృతీకరించదగినది. అంటే,
నిర్దుష్టమైనవీ సంపూర్ణమైనవీ అయిన మొదటిస్థాయి తర్క నిరూపణ వ్యవస్థలు
ఉన్నాయి; అంటే ఏ \tetoken{వాక్యాల}{!!{sentence}s} సమితి~$\Gamma$కు, ఏ
\tetoken{వాక్యం}{!!{sentence}}~$!A$కు అయినా $\Gamma \Entails !A$ అయితేనే
$\Gamma \Proves !A$ అయ్యే ధర్మం గల \tetoken{వ్యుత్పాద్యతా}{!!a{derivability}}
సంబంధం~$\Proves$ను అవి ఇస్తాయి. దీని అర్థం ముఖ్యంగా, మొదటిస్థాయి
తర్కపు చెల్లుబాటయ్యే వాక్యాలు \tetoken{గణనీయంగా
  లెక్కించదగినవి}{!!{computably
  enumerable}}. $f$ విలువలు సరిగ్గా $\Lang{L}$లోని చెల్లుబాటయ్యే
\tetoken{వాక్యాలే}{!!{sentence}s} అయ్యే గణనీయ ప్రమేయం~$f\colon \Nat \to
\Sent[L]$ ఒకటి ఉంది. ఎందుకంటే \tetoken{వ్యుత్పత్తులను}{!!{derivation}s} లెక్కించవచ్చు;
వాటిలో ఒక్క \tetoken{వాక్యాన్ని}{!!{sentence}} \tetoken{నిగమించేవాటిని}{!!{derive}} ఆ
\tetoken{వాక్యానికి}{!!{sentence}} నిర్దేశించవచ్చు. మొదటిస్థాయి తర్కం కంటే
ద్వితీయ-స్థాయి తర్కానికి వ్యక్తీకరణశక్తి ఎక్కువ; అందువల్ల సాధారణంగా
దాని చెల్లుబాటయ్యే వాక్యాలను పట్టుకోవడం మరింత సంక్లిష్టం. నిజానికి,
ద్వితీయ-స్థాయి తర్కం అనిర్ణయనీయమైనదే కాక, దాని చెల్లుబాటయ్యే వాక్యాలు
\tetoken{గణనీయంగా
  లెక్కించదగినవి}{!!{computably
  enumerable}} కూడా కావని చూపుతాం. దీని వల్ల ద్వితీయ-స్థాయి తర్కానికి
నిర్దుష్టమైన, సంపూర్ణమైన నిరూపణ వ్యవస్థ ఏదీ ఉండదు (అయితే నిర్దుష్టమైన,
కాని అసంపూర్ణమైన నిరూపణ వ్యవస్థలు ఉన్నాయి; అవి నిజానికి ముఖ్యమైన
పరిశోధనా విషయాలు).

మొదటిస్థాయి తర్కానికి మరో రెండు ధర్మాలు కూడా ఉన్నాయి: అది సంహతమైనది
(ఒక \tetoken{వాక్యాల}{!!{sentence}s} సమితి~$\Gamma$లోని ప్రతి పరిమిత ఉపసమితి
సంతృప్తిపరచదగినదైతే, $\Gamma$ స్వయంగా సంతృప్తిపరచదగినది); దానికి
లొవెన్‌హైమ్--స్కోలెమ్ సిద్ధాంతం వర్తిస్తుంది ($\Gamma$కు ఒక అనంత
నమూనా ఉంటే, దానికి \tetoken{ఒక అనంతంగా లెక్కించదగిన}{!!a{denumerable}} నమూనా
ఉంటుంది). ఈ రెండు ఫలితాలూ ద్వితీయ-స్థాయి తర్కానికి విఫలమవుతాయి.
ద్వితీయ-స్థాయి తర్కం \tetoken{వ్యక్తి క్షేత్రాల}{!!{domain}s} పరిమాణం గురించిన,
మొదటిస్థాయి తర్కం వ్యక్తం చేయలేని వాస్తవాలను వ్యక్తం చేయగలగడమే
మళ్లీ దీనికి కారణం.

\end{document}
